\chapter{LITERATURE REVIEW}


\section{2.1 Machine learning in solar photovoltaic systems}
Solar photovoltaic (PV) systems are highly dependent on environmental factors such as, temperature, solar irradiance and weather variability, which makes accurate prediction and system optimisation challenging. The application of Machine Learning (ML) techniques to renewable energy, specifically solar energy forecasting, has become increasingly prominent and has been extensively studied in order to deal with the fluctuation and uncertainty of solar power generation. The literature indicates that forecasting-related applications dominate PV research, as precise prediction is necessary for grid stability and efficient energy management. ML models can learn complex relationships between environmental factors and PV output, making them more effective than traditional statistical approaches \parencite{Massidda et al._2025}. Before this, energy forecasting depended on statistical and physics-based models, which often fail to accurately represent complex nonlinear relationships between meteorological variables and power output under dynamic environmental conditions \parencite{Siddiqui et al._2020} These traditional methods are significantly reliant on high-quality, site-specific data and often exhibit poor performance when dealing with incomplete datasets, hence limiting their effectiveness in practical applications. \parencite{Yousef et al._2020}. On the other hand, ML-based approaches demonstrated enhanced proficiency in learning patterns from large-scale data, enabling more precise and adaptive forecasting of renewable energy generation \parencite{Fernando et al._}. Nonetheless, despite their robust predictive performance, many of these models operate as "black boxes", lacking transparency and interpretability, which poses significant challenges for their application in critical energy systems \parencite{Yousef et al._2020}. Black boxes are the models where the mathematical formulations of the relationship between input and output are not fully revealed. To overcome this limitation, explainable artificial intelligence (XAI) methods have been progressively integrated into renewable energy forecasting models to provide insights into model decision-making processes (Fernando et al.). These methods enable the identification of key environment factors influencing energy generation and improve model transparency while maintaining high predictive accuracy (Yousef et al.). In general, the integration of ML and XAI techniques supports the development of intelligent and interpretable forecasting systems that improve decision-making in renewable energy management (Fernando et al; Siddiqui et al.). 

\section{2.2 Overview of traditional machine learning for solar photovoltaic}
A variety of traditional machine learning algorithms have been implemented for solar photovoltaic systems and analysis due to their effectiveness in handling structured data \parencite{Massidda et al._ 2025}. The traditional machine learning models are able to capture non-linear relationships between input data and PV output. Ensemble techniques like bagging and boosting further enhance performance by combining multiple models, resulting in better aacuracy and robustness. Furthermore, traditional ML models use less data and are more computationally efficient compared to deep learning techniques. This makes them suitable for real-world scenarios, especially in resource-constrained environment. Another important feature is their partial interpretability. For example, tree-based models are more transparent than black-box models because they provide feature importance measurements that aid in understanding the contribution of input variables \parencite{De Silva et al._2024}. This section provides an overview of some commonly used traditional ML algorithms found in the reviewed literature.

\subsection{2.2.1 Extreme Gradient Boosting (XGBoost) }
Extreme Gradient Boosting (XGBoost) is an optimized gradient boosting algorithm that builds a stronger, high-performance model by combining several weak models. Boosting is a method of using decision tree as base learners and building them sequentially, so each tree corrects errors from the previous one. XGBoost enables parameter customization to optimize performance on particular issues and has parallel processing for faster training on large datasets. Its capability to handle complex and high dimensional data makes it very effective, which is an invaluable advantage in applications like solar photovoltaic (PV) systems. XGBoost has demonstrated exceptional performance in classification tasks in several studies, particularly in predicting energy output in PV systems and distinguishing between various types of partial discharge signals. For example, \parencite{Kumar et al._2020} discovered that XGBoost provided better accuracy in classifying defects in PV systems compared to other traditional machines learning algorithms. In addition, according to \parencite{Zhang et al._2020}, XGBoost consistently outperformed more straightforward models like k-Nearest Neighbors and Decision Trees in fault detection and energy forecasting. XGBoost is a preferred option in solar PV applications due to its robustness in multi-class classification problems and its ability to handle large-scale, especially where accurate predictions are crucial for efficiency and maintenance \parencite{Shen et al._2020}. Furthermore, XGBoost’s performance is often enhanced when paired with advanced feature selection techniques, ensuring its efficiency across a range of datasets and operational conditions. Traditional models such as decision tree and random forest are straightforward and easy to interpret but may not be accurate when dealing with complex data. This highlights the advantage of XGBoost in more challenging scenarios.

\subsection{2.2.2 Support Vector Machines (SVM)}
Support vector machines (SVM) is a supervised machine learning approach used for problems involving regression and classification. It tries to find the hyperplane, which is the best boundary that separates different classes in the data. SVM’s main goal is to maximize the margin between the two classes. SVM is useful for binary classification. Support vector machines (SVM) have proven to be a powerful and widely used machine learning technique in solar photovoltaic systems for classifying PV performance data. For example, identifying faults or predicting the system’s energy production based on environmental conditions. According to \parencite{Kumar et al._2020}, SVM outperforms other traditional machine learning models like Decision Trees and k-Nearest Neighbors (K-NN) in terms of classification accuracy when it comes to distinguishing between different operational states of PV systems. \parencite{Siddiqui et al._2020} added that SVM is robust even when used on relatively smaller datasets, which is a typical limitation in PV applications where collecting vast amounts of labelled data might be difficult. One of the key strengths of SVM in PV systems is its capability to handle high-dimensional data, such as environmental and time variable, while maintaining high accuracy in system performance prediction and anomaly detection. \parencite{Shen et al._2020} pointed out that SVM can effectively capture complex patterns in PV data, which are frequently influenced by weather conditions, irradiance levels, and panel orientation. This is due to the SVM’s ability to handle non-linear relationships through kernel functions like the radial basis function (RBF) kernel. Therefore, SVM remains a preferred choice for many solar PV applications due to its robustness, accuracy, and flexibility in modelling complex relationships in energy generation and fault detection.

\section{2.3 Feature engineering in solar photovoltaic systems}
Feature engineering plays a critical role in improving the performance of ML models in PV systems. It involves selecting, transforming, and creating new features from raw data to improve model accuracy and prediction power. In context of solar PV systems, various types of features arise from time-series data that show daily or seasonal variations in PV performance, as well as from operational data like temperature, irradiance, and energy output. Statistical measures, including the mean, variance, and skewness are commonly used to describe the general behaviour of the system \parencite{Dinakar et al._2021}. These features can help distinguish between normal operation and potential faults or inefficiencies in the PV system. 

\subsection{2.3.1 Time-domain features}
Time-domain features are important in the analysis and monitoring of solar photovoltaic (PV) systems as they provide insights into the instantaneous behaviour of the system. These features are derived directly from the raw time-series data, which capture important aspect of the signals related to energy production, fault occurrence, and system performance. For example, peak amplitude is a commonly used time-domain feature that represents the maximum value of the signal during a specific period. It is directly associated with the intensity of energy generation or potential faults in the system \parencite{Kumar et al._2020}. Another crucial time-domain feature is pulse duration, which measures the duration of a discharge or power fluctuation. This feature is indicative of the behaviour of defects. This is because various fault types usually exhibit in different pulse durations, including internal voids or surface discharge \parencite{Shen et al._2020}. Furthermore, the rise time of the signal reflects the dynamics of energy changes in PV systems. The rise time of the signal measures how quickly the signal reaches its peak value from 10% to 90%. This is particularly in relation to the speed of ionization processes in energy conversion \parencite{Hayawi et al._2021}. The rise time in PV systems can be very useful in distinguishing between different operational conditions, such as normal fluctuations in power output and those caused by faults or environmental interference \parencite{Zhang et al._2020}. Another crucial time-domain aspect is pulse count. This refers to the number of discharges or power fluctuations that occur within a single operating cycle. Higher pulse counts are crucial for fault detection and predictive maintenance because they may reveal more severe or active insulation defects \parencite{Kumar et al._2020}. These time-domain features are essential for modelling and diagnosing the performance of solar PV systems. This is because they assist distinguish bot normal and abnormal operating states. This allows for the early detection of any issues and facilitating better maintenance plans. 
\subsection{2.3.2 Frequency-domain features}
In solar photovoltaic (PV) systems, frequency-domain features are essential for identifying patterns that are not easily observable in the time-domain. These features are obtained by transforming time-series data into the frequency domain, usually using techniques like wavelet or Fourier transforms. The frequency-domain analysis allows for the identification of periodic components, which can offer important insight into the behaviour of PV systems, especially when handling issues like shading, temperature variations, or system degradation. For example, anomalies or changes in the energy output that occur at particular frequency can be identified using frequency-domain feature, which may point to system inefficiencies or fault conditions \parencite{Shen et al._2020}. In PV systems, the spectral density of the signal can reveal underlying periodic behaviours that affect energy production. The spectral density of the signal represents the distribution of power across different frequencies. According to \parencite{Kumar et al._2020}, these features are crucial for fault identification because they allow for the separation of normal operational fluctuations from those caused by defects or external disturbances. An additional crucial stage in fault identification is the analysis of frequency bands and their energy distribution, since certain defects might lead to distinct spectral patterns \parencite{Shen et al._2020}. Furthermore, frequency-domain features can be useful in identifying resonant frequencies where PV system components may experience amplified energy oscillations and perhaps leading to malfunction or failure. As noted by \parencite{Zhang et al._2020}, frequency-domain analysis can help distinguish between normal system behaviour and early signs of mechanical or electrical degradation. This provides deeper insights into the overall condition and performance of the PV system, including its components such as solar panels, inverters, and electrical connections. Overall, frequency-domain features improve the ability of machine learning to analyze PV system performance, making them crucial component for improving fault detection, energy forecast, and system optimization.
\section{2.4 Dataset diversity and model generalisability}
The performance and reliability of ML models in PV systems are significantly affected by dataset diversity. Numerous research relies on datasets collected from specific geographical locations, which restricts the ability of models to generalise across different environments \parencite{Massidda et al._2025}. When applied to new locations, model performance can be greatly affected by variations in climate, weather patterns and system configurations. As a result, models that were trained on a single dataset might not perform well in real-world scenarios. The literature highlights the significance of using a variety of datasets and conducting cross-site validation in order to improve generalisability. Public datasets such as the National Solar Radiation Database (NSRDB) are often used to enhance reproducibility and consistency \parencite{De Silva et al._2024}. 
\section{2.5 Comparative studies and benchmarking in machine learning for solar photovoltaic systems}
Comparative studies are essential for evaluating the performance of different ML models in PV systems. The research shows that ensemble methods such as Random Forest and gradient boosting often outperform single models due to their ability to reduce variance and improve robustness \parencite{Massidda et al._2025}. However, the lack of standardised datasets and evaluation metrics makes benchmarking across studies remains challenging. Direct comparison is challenging due to the varying forecasting horizons, input features and validation methods used in different studies. Common evaluation metrics include Root Mean Square Error (RMSE), Mean Absolute Error (MAE) and Mean Absolute percentage Error (MAPE). These metrics may not accurately capture model performance under real-world scenarios despite the fact that they provide useful insight \parencite{De Silva et al._2024}. According to the research, more thorough benchmarking frameworks are required, including cross-site validation and testing under noisy or missing data situations. 
\section{2.6 Problems and limitations of traditional machine learning}
Traditional ML models encounter several challenges in PV application despite of their advantages. Data quality is a significant issue since real-world PV data has sensor errors, noise and missing values. If these issues are not appropriately resolved, it can significantly degrade model performance. Another limitation is lack of generalisation, where models trained on one dataset may not perform well in different environmental conditions. Class imbalance is also a concern, especially in fault detection tasks where abnormal conditions are uncommon. This may lead to biased models that are unable to detect faults accurately. Furthermore, it is challenging to compare findings across studies due to the lack of standardised evaluation methods. This problem is due to the differences in datasets, forecasting timeframes and performance metrics. Finally, some traditional ML models still exhibit black-box behaviour, which restricts their interpretability and reliability.
\section{2.7 Importance of Explainable Artificial Intelligence (XAI)}
Explainable Artificial Intelligence (XAI) has become increasingly important in PV systems to overcome the limitations of ML models. XAI provides transparency by explaining how input features influence model predictions, thereby improving trust and reliability. For model interpretation, techniques like SHAP and LIME are widely used. SHAP provides both local and global explanations based on a strong theoretical framework, while LIME focuses on explaining specific predictions. In order to ensure that models rely on physically meaningful variables rather than inaccurate correlations, XAI is also crucial in feature selection and validation. This enhances both scientific reliability and model performance. Additionally, XAI offers robustness analysis and model debugging, allowing researchers to evaluate model behaviour under noisy or incomplete data conditions. This is crucial for real-world applications, where data imperfections are common.
\section{2.8 Research gaps}
Several research gaps can be identified from the existing literature. First, many studies fail to sufficiently address real-world data issues such as noise, missing values and sensor errors. Second, there is limited focus on dataset diversity and cross-site validation, which affects model generalisability. Third, it is still challenging to compare various machine learning approaches due to inconsistent benchmarking practices. Lastly, the integration of XAI techniques is still not thoroughly investigated although traditional ML models offer some level of interpretability. These gaps highlight the need for more robust, interpretability and generalisable ML models for PV applications.
\subsection{Planning Your Literature Search}

\subsection{Searching for Relevant Literature}
Develop a search strategy using relevant keywords, phrases, and synonyms. Apply appropriate filters and Boolean operators (AND, OR, NOT) to refine your search results. Record your search strategies and results, which is essential for the reproducibility of your research.

\subsection{Screening and Selecting Literature}
Screen titles and abstracts for relevance to your research questions. Obtain and read full texts for selected sources. Use inclusion and exclusion criteria to systematically decide which articles to consider in your review.

\subsection{Analyzing and Synthesizing Information}
Organize the selected literature into thematic categories or according to methodological approaches. Analyze the findings and methodologies used in the literature to identify patterns, themes, and gaps in the research.

\subsection{Writing the Review}
Summarize the literature, linking it directly to your research questions and objectives. Discuss the significance of findings in relation to previous studies. Highlight any controversies, inconsistencies, and gaps in the literature. Present a critical analysis of the collected data, providing your interpretations and insights.




\begin{table}[h]
\centering
\caption{Comparison of Different Techniques}
\label{tab:technique_comparison}
\begin{tabular}{|c|m{6cm}|m{6cm}|}
\hline
\textbf{No.} & \textbf{Brief Note} & \textbf{Additional Details} \\ \hline
1 & Technique A: This technique is known for its speed and efficiency in processing large datasets. & Suitable for real-time applications. It can handle streaming data effectively but may require significant computational resources. \\ \hline
2 & Technique B: Often used for its high accuracy in classification tasks, especially in machine learning. & Requires extensive training data and is computationally expensive. It excels in scenarios where precision is critical. \\ \hline
3 & Technique C: Valued for its flexibility and adaptability across various domains. & While versatile, it may not provide the best performance in terms of speed or accuracy in specific applications. \\ \hline
\end{tabular}
\end{table}



\begin{table}[]
	\centering
	\caption{Reported results of the ternary sleepiness detection method in previous studies. }
	\label{tab:table3levelCompareOtherMinimal}
	\begin{tabular}{lllllll}
		\toprule[\heavyrulewidth] \toprule[\heavyrulewidth]
		No& Study   & Parameter(s) & Classifier   \\
		\midrule[\heavyrulewidth]
		1&\cite{Barua_2019} & 	EEG, EOG, CF& SVM \\
		2&\cite{Chai2019} & SWA & SVM \\
		3&\cite{Wang2016} & PERCLOS, SWA & 	MLO\\
		4&\cite{Zhao2015} & FE & NB \\
		5&\cite{Zhao2018} & FE & NB \\
		6&\cite{Zilin2017} & PERCLOS, EM & 	ANN\\
		7&\cite{Picot_2012} & 		EEG, EOG & 		CDR \\
		\hline	
	\end{tabular}
\end{table}

\section{Research Gap}

In the research gap section, it's important to find areas where questions still need answers or where current theories and methods don't fully solve the problems. This involves combining knowledge from previous studies you've reviewed, critically looking at their findings and methods, and spotting areas that haven't been fully explored yet. Doing this should help you see where you can make important improvements in the field, suggest new ways of researching, or address complex issues that haven't been dealt with thoroughly in past research.

\section{Chapter Summary}

A chapter summary in the literature review section is a brief overview of the key points covered in that chapter. It helps to consolidate the main findings, theories, and discussions presented. In this summary, you should clearly outline the most important ideas and how they relate to your research topic. This includes summarizing the debates, methodologies, and conclusions from various sources you've studied. The purpose of this summary is to give a clear and concise recap that helps readers understand the background and context of your study, and to set the stage for presenting the research gaps and questions your study aims to address.
